\section{Scope, Method, and Qualifications}

\subsection{Reader, question, and method}

This is a discursive theological article for a serious general reader, able to follow careful distinctions but not presumed to know scholastic terminology. It is the fourth and closing volume of a four-part series on the last things; each volume is a complete study of its own subject, and this one neither summarises nor depends on the other three. Where its subject touches theirs---the particular judgment presupposed by \emph{Benedictus Deus}, the purification the same constitution presupposes, the possibility of definitive refusal---it states only what its own sources establish and refers to the companion volumes by their subjects.

The governing question is what may be said about heaven on evidence that can be produced, and its claim is twofold. First: what the Catholic Church has defined is a promise of an immediate, unmediated, personal seeing of God himself, given to creatures who cannot achieve it, differing in degree with charity, entailing the resurrection of the very bodies we now bear, and enjoyed in company---and this defined core is both smaller and more extreme than the picture that has grown over it. Second: the great majority of what an ordinary Catholic believes about heaven is theological speculation, and a substantial part of that was flagged as conjecture by the authors who produced it and un-flagged by their readers. The second claim is project synthesis and is stated with its counterargument in the synopsis; the first is a report of the definitions and their limits.

The method is one discipline applied at every consequential claim: name the register. Defined doctrine, condemned error, catechetical or conciliar teaching, Scripture, patristic witness, school theology of either tradition, reported reception, and original synthesis are labelled where they occur and again in the synopsis. Weight is not inferred from eloquence, antiquity, or frequency of repetition. Where a claim rests on a witness the article could not reach at its own locus, the ceiling is stated in the note and in the source audit rather than smoothed over.

The body proceeds in the order the evidence has rather than in a topical inventory: the definitions first, because they are the fixed points; then the question the definitions raise and do not answer, worked in the \emph{Summa} at the text; then the Eastern alternative in its own authors, placed before the remaining topics because it is a rival account of the same central act and not an appendix to it; then the body, the degrees, and the company, which are the three places where confident detail and actual warrant diverge most sharply; then an audit of the popular picture against bounded searches of the tracked Scripture text; then the close.

\begin{threestable}{Source class}{Function in the article}{Governing boundary}
Defining acts & \emph{Benedictus Deus} (1336) supplies the central definition, read at its Latin phrase by phrase; Florence's \emph{Laetentur caeli} (1439) adds \latin{sicuti est} and the degrees; Lateran IV's \emph{Firmiter} (1215) the resurrection clause; Trent session VI canon 32 the increase of glory; Vienne's \emph{Ad nostrum qui} (1312) the condemned denial of the light of glory. & Quoted in Latin from the received text of a standard collection in a dated web presentation that does not identify the printing it transcribes; that is the stated verification ceiling. English renderings are labelled working glosses except where the Catechism's official English is quoted. Dogmatic force is confined to what the texts state.\\
Catechetical and conciliar teaching & CCC 1000, 1023--1029, 1045, 1048 on heaven, the vision, communion, and the limits of knowledge; \emph{Lumen gentium} 49--51 on the communion of saints; \emph{Unitatis redintegratio} 17 on Eastern and Western theological expression. & Cited at its own authority, which is authoritative and not definitive; brief excerpts of copyrighted official texts, quoted with attribution. Catechetical language is not treated as technical definition.\\
Scripture & The texts denying vision (Ex 33:20; Jn 1:18; 1 Tim 6:16) and promising it (Matt 5:8; 1 Cor 13:12; 1 Jn 3:2; Apoc 22:4); the resurrection chapter (1 Cor 15); the Apocalypse's final vision; and the bounded searches behind the audit of the popular picture. & Quoted from the public-domain Douay--Rheims (Challoner) in the source library's tracked per-verse text of the Project Gutenberg electronic edition, at Vulgate numbering, on 2026-07-26. Readings are textual and modest; no independent exegetical apparatus is claimed; negative results are bounded to a literal English search in one edition.\\
Latin school theology & ST I, q.~12 governs the analysis of the vision, the light of glory, comprehension, and degrees; I-II, q.~4 the body and the fellowship of friends; II-II, q.~26, a.~13 the order of charity in heaven; the Supplement qq.~82--85 and 96 the properties of the glorified body and the aureoles. & Adopted as the article's principal Western framework where Catholic theology permits schools; not asserted as dogma; each locus used at its exact argumentative level. The Supplement is identified at every use as compiled after Aquinas's death from his earlier commentary.\\
Eastern tradition & Basil (Letter 234), Chrysostom (\emph{Hom.\ on John} 15), Gregory of Nyssa (\emph{On the Soul and the Resurrection}), John of Damascus (\emph{De fide orth.} I.4) in public-domain translations; Gregory Palamas's \emph{Triads} in an identified modern translation. & Stated at full strength in its own authors, not as a foil. The Palamas quotations were read in an optical-character transcription of a scanned copy of a copyrighted 1983 translation, without the Greek and without collation; the conciliar acts of 1341, 1347, and 1351 were not reached at all, and their outcome is reported from a named modern introduction.\\
Reference works & The 1907--1912 \emph{Catholic Encyclopedia} for the sequence of the 1330s controversy, for Benedict XII's preparation of the constitution, and for the compilation of the \emph{Summa}'s Supplement. & Public-domain historical witnesses of their date, used for chronology and bibliography, never for doctrinal authority. Their judgments of probability (for example on the Supplement's compiler) are reproduced as probabilities.\\
Project synthesis & The three-register discipline; the judgment on whether the Latin and Palamite doctrines formally contradict; the judgment on whether the communion of saints is constitutive; the reading of the Apocalypse's imagery as systematically privative around one positive clause; the estimate of how much of the received picture is speculation; the synopsis table and all final formulations. & Original editorial synthesis argued from the checked sources, attributed to no source, creating no doctrine, and resolving nothing the Church has left open. The two labelled synthesis blocks state their counterargument at full strength and name what finding would change them.\\
\end{threestable}

\subsection{Included and excluded scope}

Included are: the text and construction of \emph{Benedictus Deus} and the controversy that produced it; the additions made at Florence and Lateran IV, with Trent's canon on the increase of glory and Vienne's condemned proposition; the scriptural texts that deny and that promise the vision of God; the Thomistic account of the vision, the light of glory, and non-comprehension, read in ST I, q.~12; the Eastern essence-and-energies doctrine in Basil, Nyssa, Chrysostom, the Damascene, and Palamas, with its fourteenth-century reception reported at a stated ceiling; the defined and undefined content of the resurrection of the body, with Augustine's speculations and disclaimers and the schools' four gifts; the degrees of beatitude, their measure in charity, and the ways equality and difference are held together; the communion of saints as related to the beatifying act; and a bounded audit of the popular imagery against the tracked Scripture text.

Excluded are: death, particular judgment, purgatory, the general judgment, and hell, which belong to the companion volumes and are touched here only where a source under examination mentions them; the theology of grace, merit, justification, predestination, and assurance beyond the single canon of Trent cited for the increase of glory; Christology and soteriology, presupposed throughout and nowhere developed; the sacramental economy and the Eucharist as pledge of glory; the intra-Thomist controversy over the natural desire for the supernatural, which is named and not adjudicated; the theology of the angels beyond what the sources quoted say; private revelation, mystical or near-death testimony, and any phenomenology of religious experience; the history and iconography of Christian art, which the audit of the popular picture touches only to say that a convention is a convention; the reception history of \emph{Laetentur caeli} in the East; and any judgment about the state, culpability, degree of glory, or destiny of any person.

\subsection{Material qualifications}

\begin{enumerate}[leftmargin=1.45em,itemsep=0.25em]
\item \latin{Beatitudo}, rendered ``beatitude'' or ``happiness,'' names the complete and final good of a rational creature, not a mood; the English Dominican \emph{Summa} translation's ``Happiness'' carries that sense wherever it is quoted.
\item Every Latin text of a magisterial act quoted here was read in a dated web presentation of a standard collection that does not identify the printed edition it transcribes. No printed edition was collated. This is the article's verification ceiling for all conciliar and papal wording, and it is a real one: a collation could in principle change a word.
\item The English renderings of Latin acts are the article's own labelled working glosses, governed by the Latin, except where the Catechism's official English is quoted as such. They are not approved translations and must not be used as if they were.
\item \emph{Benedictus Deus} defines the state of purified souls before the resurrection. Its use here is confined to what it states; no position is taken on eschatological questions outside it, and the constitution's final paragraph on those who die in mortal sin is noted and left to the volume that owns it.
\item The exclusion in the definition is of a created \emph{object} of vision. The compatibility of the Thomistic light of glory with that exclusion depends on the scholastic distinction between a medium \latin{in quo} and a medium \latin{quo}; the article states the distinction, shows that Aquinas draws it explicitly, and does not claim that the definition canonises his account.
\item The Council of Vienne's condemned proposition is a compound. Its condemnation establishes that natural beatitude coupled with the denial of any elevating light is not an open option; it does not define the Thomistic explanation of the light of glory.
\item How much the argument from natural desire (ST I, q.~12, a.~1) proves has been disputed among Thomists for centuries. The article names the dispute, does not adjudicate it, uses the argument only as an explanation of a datum already held on other grounds, and nowhere treats the desire as an exigency binding God.
\item The Palamas material is quoted from Nicholas Gendle's English translation in the 1983 Paulist Press edition, read in an optical-character transcription of a scanned copy hosted at the Internet Archive. The Greek was not consulted; no printed copy was collated; the transcription contains visible uncorrected scanning errors elsewhere in the file. Quotations are short and attributed. This is the article's stated ceiling for every Palamas sentence.
\item The acts of the Byzantine councils of 1341, 1347, and 1351, the Hagioritic Tome, and the Synodal Tome were not reached in any edition. Their existence and outcome are reported from a named twentieth-century scholarly introduction, and no claim whatever is made about their wording, canons, or anathemas. This is the largest single gap in the article's evidence and it bears directly on the synthesis about contradiction, which is accordingly held at low confidence.
\item The judgment that \latin{essentia} and \latin{ousia} are probably not the same concept is a proposal about where to look, not a finding, and is stated with the counterargument at full strength and with the specific finding that would overturn it. It resolves nothing between the Churches and has no ecumenical standing.
\item The judgment that the communion of saints is constitutive of the state but not of the beatifying act is likewise labelled synthesis and is held at moderate confidence. \emph{Lumen gentium} 51's boundary---that communion with the saints must not weaken the worship due to God---is stated beside it and is not this article's qualification but the Council's.
\item The Supplement to the \emph{Summa theologiae} was assembled after Aquinas's death from his commentary on the \emph{Sentences}; the compiler's identity is not settled, and the older reference work cited gives the attribution to Reginald of Piperno as probable, not certain. The material is authentic Thomas in substance and early Thomas in date, transposed by another hand.
\item Augustine's speculations on the risen body are quoted from the public-domain Dods translation of 1871 in the source library's tracked transcription of that printing. His two disclaimers are quoted in full because the article's use of the material depends on them; nothing is attributed to him that he did not say, and the tentativeness of his register is reported as he marked it.
\item The negative results reported in the audit are bounded literal searches for English strings within single verses of one electronic edition of one translation, run on 2026-07-26. They cannot exclude synonyms, paraphrase, inflection, a different translation's vocabulary, or a sense expressed without the word. ``Not found'' means exactly that and never ``false.''
\item The tracked Douay--Rheims text used here preserves Vulgate chapter-and-verse numbering, which differs from a modern Bible's at several loci cited, most consequentially at 1 Thess 4:16. It also differs from the 1899 American edition in 1859 verses, collated in a table registered with the edition; where a divergence touches a verse quoted here (Job 19:26, Apoc 22:2) the note says so, and no claim depends on the difference.
\item The estimate that ``the great majority'' of the received popular picture is speculation is a qualitative judgment, not a measurement. No corpus of devotional literature was sampled and no count was made; the judgment rests on the sorting performed in the body, and a reader who thinks the proportion smaller is entitled to the same evidence.
\item No pastoral, canonical, or liturgical advice is given. Nothing here diagnoses any person's state or destiny, and nothing here is offered as consolation in the place of a competent pastor.
\end{enumerate}

\subsection{Notes, translations, rights, currentness, and review}

The numbered Notes carry exact citations, edition and translation identifications, verification dates, search bounds, and the qualifications that would interrupt the argument; no indispensable premise lives only in a note. Scripture is quoted from the public-domain Douay--Rheims (Challoner revision) in the source library's tracked per-verse text of the Project Gutenberg electronic edition. Aquinas is quoted from the public-domain English Dominican translation, with the Latin of I, q.~12 checked at Corpus Thomisticum; the New Advent presentation inserts spaces before punctuation at hyperlinked words, and quotations here close those spaces and change nothing else. Basil, Chrysostom, Gregory of Nyssa, and John of Damascus are quoted from the public-domain Nicene and Post-Nicene Fathers translations; Augustine's \emph{De civitate Dei} from the public-domain Dods translation of 1871, checked in the source library's tracked transcription of that printing. Latin conciliar and papal texts are quoted from the received text of the \emph{Enchiridion symbolorum} in a dated public web presentation, with English renderings labelled as working glosses. The Catechism, \emph{Lumen gentium}, and \emph{Unitatis redintegratio} are quoted briefly from the official Vatican web texts with attribution. Gregory Palamas is quoted in short clauses from an identified copyrighted modern translation, with attribution and with the transcription ceiling stated. Official texts and third-party translations remain outside the project's CC~BY~4.0 grant; project-created prose, organisation, and synthesis are project content.

All online witnesses were checked on 2026-07-26; no mutable law or discipline is relied on, and no as-of date governs any conclusion. This revision received internal argumentative, source-consistency, quotation, rights, and production review by the authoring agent. Independent scriptural, patristic, Thomistic, Byzantine, dogmatic, theological, and ecclesiastical review is outstanding; no imprimatur, nihil obstat, or ecclesiastical approval is claimed, and internal checking is not independent review.
